\documentclass[12pt]{article}

\usepackage[a4paper, margin=1in]{geometry}
\usepackage{setspace}
\usepackage{amsmath}
\usepackage{graphicx}
\usepackage{hyperref}
\usepackage{natbib}
\usepackage{titlesec}

\setstretch{1.25}

\title{\textbf{Machine Learning for Time Delay and Light Curve Analysis in Astroinformatics}}
\author{Jorge Alberto Suárez Saldaña \\
Machine Learning Course \\
Advisor: Professor Cuevas Tello Juan Carlos}
\date{\today}

\begin{document}

\maketitle
\newpage

\section{Introduction}

Astroinformatics combines astronomy, statistics, and computer science to extract knowledge from large-scale astronomical datasets. One of the most important data structures in time-domain astronomy is the \textbf{light curve}, which represents the variation of a celestial object's brightness as a function of time.

Light curves are essential for detecting and analyzing:
\begin{itemize}
    \item Exoplanet transits
    \item Stellar rotation and variability
    \item Supernovae
    \item Gravitational lensing phenomena
\end{itemize}

Closely related is the concept of \textbf{time delay}, which refers to measurable shifts between light signals caused by gravitational lensing or geometric differences in light travel paths. Time-delay measurements are fundamental in cosmology and allow estimation of distances and even cosmological parameters.

With modern telescopes such as \textit{Kepler}, \textit{TESS}, and the upcoming \textit{Vera Rubin Observatory (LSST)}, astronomy has entered a data-intensive era. Traditional analysis methods struggle to scale efficiently, creating a strong motivation for Machine Learning-based approaches.

\section{Problem Definition and Motivation}

The core problem addressed in this project is:

\begin{quote}
How can Machine Learning methods improve the detection, characterization, and time-delay estimation of astronomical events using light curve data?
\end{quote}

Light curves are typically irregularly sampled time-series data affected by:
\begin{itemize}
    \item Instrumental noise
    \item Missing observations
    \item Atmospheric interference
    \item Earth's orbital motion (requiring Barycentric corrections)
\end{itemize}

Classical astronomy software relies heavily on:
\begin{itemize}
    \item Lomb-Scargle periodograms
    \item Fourier transforms
    \item Cross-correlation methods
    \item Template fitting
\end{itemize}

However, these approaches:
\begin{itemize}
    \item Struggle with irregular sampling
    \item Require manual parameter tuning
    \item Become computationally expensive at large scale
\end{itemize}

Modern sky surveys generate terabytes of data per night. Human inspection is no longer feasible. Machine Learning offers pattern recognition capabilities that can automatically detect subtle variations across massive datasets.

This motivates the application of AI methods to:
\begin{itemize}
    \item Detect periodic behavior
    \item Estimate transit parameters
    \item Measure gravitational lensing time delays
    \item Classify variable astronomical sources
\end{itemize}

\section{State of the Art}

\subsection{Classical Approaches}

Traditional methods include:

\begin{itemize}
    \item Lomb-Scargle periodograms for irregular time-series
    \item Cross-correlation for time-delay estimation
    \item Bayesian model fitting
    \item Template matching
\end{itemize}

These methods are mathematically well-founded but computationally intensive when applied to millions of light curves.

\subsection{Machine Learning Approaches}

Recent research shows strong adoption of ML methods:

\begin{itemize}
    \item Gradient Boosted Trees for photometric classification
    \item Random Forests for transient detection
    \item Convolutional Neural Networks (CNNs) for transit detection
    \item Recurrent Neural Networks (RNNs, LSTM) for sequential modeling
    \item Gaussian Processes (GPs) for uncertainty modeling
\end{itemize}

Deep learning allows:
\begin{itemize}
    \item Automatic feature extraction
    \item Robust handling of noisy signals
    \item Scalable classification of millions of sources
\end{itemize}

Recent works demonstrate that ML models achieve competitive or superior performance compared to classical methods, particularly when dealing with large and irregular datasets.

\subsection{Contrast: Classical vs Machine Learning}

\begin{center}
\begin{tabular}{|p{3cm}|p{5cm}|p{5cm}|}
\hline
Aspect & Classical Methods & Machine Learning \\
\hline
Feature Design & Manual & Automatic / Learned \\
\hline
Scalability & Limited & Highly scalable \\
\hline
Irregular Data & Sensitive & Robust (especially RNN/GP models) \\
\hline
Interpretability & High & Moderate \\
\hline
Computation After Training & High per case & Fast inference \\
\hline
\end{tabular}
\end{center}

Machine Learning shifts the paradigm from model-fitting to data-driven pattern discovery.

\section{Methodology}

The proposed methodology consists of:

\subsection{Data Preprocessing}
\begin{itemize}
    \item Convert timestamps to Barycentric Julian Date
    \item Normalize flux values
    \item Remove outliers
    \item Handle missing data
\end{itemize}

\subsection{Feature-Based Approach}
\begin{itemize}
    \item Statistical descriptors (mean, variance, skewness)
    \item Periodogram peak extraction
    \item Transit depth and duration estimation
\end{itemize}

Apply:
\begin{itemize}
    \item Random Forest
    \item XGBoost
    \item Support Vector Machines
\end{itemize}

\subsection{Deep Learning Approach}
\begin{itemize}
    \item 1D Convolutional Neural Networks
    \item LSTM / GRU models
    \item Transformer-based time-series models
\end{itemize}

\subsection{Gaussian Process Modeling}
Use scalable GP variants for:
\begin{itemize}
    \item Gap filling
    \item Uncertainty estimation
    \item Forecasting
\end{itemize}

\subsection{Evaluation Metrics}
\begin{itemize}
    \item Accuracy
    \item Precision/Recall
    \item Log-loss
    \item RMSE for regression tasks
\end{itemize}

Comparison will be made against classical baselines.

\section{Datasets}

\subsection{Public Datasets}

\begin{itemize}
    \item Kepler Mission light curves
    \item TESS photometric data
    \item LSST / PLAsTiCC simulated dataset
\end{itemize}

Typical structure:

\begin{itemize}
    \item Time (Julian Date)
    \item Flux (brightness)
    \item Flux error
    \item Object metadata (RA, DEC, redshift, etc.)
\end{itemize}

\subsection{Local Observatory Data}

If available, the university observatory dataset may include:

\begin{itemize}
    \item Time-stamped photometric measurements
    \item Telescope metadata
    \item Observational conditions
\end{itemize}

The pipeline will be designed to flexibly ingest FITS or CSV formatted light curves.

\section{Modern Digital Tools}

Modern astroinformatics workflows use:

\begin{itemize}
    \item Python ecosystem (Astropy, Lightkurve, AstroML)
    \item Scikit-learn
    \item TensorFlow / PyTorch
    \item Gaussian Process libraries
    \item Fedora Astronomy Spin (pre-configured astronomy software distribution)
\end{itemize}

These tools integrate classical astronomy processing with ML experimentation environments.

\section{Conclusion}

The transition from classical astronomical signal-processing to Machine Learning represents a paradigm shift in astroinformatics.

While classical methods remain foundational, Machine Learning enables:

\begin{itemize}
    \item Scalable processing of massive surveys
    \item Robust pattern recognition in noisy data
    \item Automated feature extraction
    \item Faster inference after training
\end{itemize}

Given the exponential growth of astronomical data, ML is not merely an enhancement—it is becoming a necessity for modern time-domain astronomy.

\newpage
\section{References}

\begin{enumerate}
\item Gabruseva et al. (2019). Photometric light curves classification with machine learning.
\item Jiménez-Vicente et al. (2025). Analysing quasar microlensing light curves with machine learning.
\item LSST DESC (2018). Photometric LSST Astronomical Time-Series Classification Challenge.
\item Priest et al. (2024). Machine learning tool fills in satellite light curves.
\item Eastman (2010). Barycentric Julian Date correction explanation.
\end{enumerate}

\end{document}
